% Thème Git : barre orange Git avec petit graphe de commits décoratif, chaque colonne est
% une branche (rail coloré, un commit sur chaque bloc), mini-graphes de commits en TikZ.
% Utilisé par Git/git_cheat.tex.
%   colonne[branche][couleur] · graphe (tikzpicture centrée) · \oui · \non

% couleurs Git (logo) : orange F05032, brun foncé 3E2C00
\definecolor{git-orange}{HTML}{F05032}
\definecolor{git-dark}{HTML}{3E2C00}
% catégories (liseré de l'onglet, sous-titres)
\colorlet{cat-base}{git-orange}  % zones, enregistrer
\colorlet{cat-his}{mk-yellow}    % historique, références
\colorlet{cat-br}{mk-green}      % branches, fusion
\colorlet{cat-rem}{mk-cyan}      % dépôts distants
\colorlet{cat-rw}{mk-purple}     % réécrire l'historique
\colorlet{cat-und}{mk-pink}      % annuler, réparer
\colorlet{cat}{cat-base}
\colorlet{rail}{git-orange}
% code : [1] sous-commandes · [2] options · [3] refs · [4] git
\colorlet{code-kw4}{git-orange}

% --- Barre d'en-tête : orange Git, petit graphe de commits décoratif ---
\colorlet{barre-fond}{git-orange}
\colorlet{barre-bord}{git-dark}
\colorlet{barre-sep}{white!65!git-orange}
\renewcommand\iconebarre{E702}         % logo Git
\renewcommand\tailleiconebarre{13}
\renewcommand\legende{\leg{git-orange}{git} \ \leg{mk-pink}{commande} \ \leg{mk-cyan}{--option} \
  \leg{mk-green}{ref} \ \legit{mk-grey}{\# commentaire} \ \legit{mk-purple}{«à remplir»}}
% graphe décoratif entre le sous-titre et la légende : main + une branche qui part et revient
\renewcommand\decorbarre{%
  \coordinate (g0) at ([xshift=7mm]sous.east|-0,\barrey);
  \coordinate (g1) at ([xshift=-7mm]leg.west);
  \draw[white,line width=1pt] (g0) -- (g1);
  \draw[white,line width=1pt,rounded corners=1.2mm]
    ($(g0)!0.18!(g1)$) |- ($(g0)!0.36!(g1)+(0,2.1mm)$) -- ($(g0)!0.62!(g1)+(0,2.1mm)$) -| ($(g0)!0.76!(g1)$);
  \draw[white!70!git-orange,line width=1pt,rounded corners=1.2mm]
    ($(g0)!0.46!(g1)$) |- ($(g0)!0.58!(g1)+(0,-2.1mm)$) -- ($(g0)!0.92!(g1)+(0,-2.1mm)$);
  \foreach \t in {0,0.18,0.32,0.46,0.6,0.76,0.88,1}
    \fill[white] ($(g0)!\t!(g1)$) circle (0.85mm);
  \foreach \t in {0.4,0.53,0.66}
    \filldraw[fill=git-orange,draw=white,line width=0.7pt] ($(g0)!\t!(g1)+(0,2.1mm)$) circle (0.75mm);
  \foreach \t in {0.7,0.82,0.92}
    \filldraw[fill=git-orange,draw=white!70!git-orange,line width=0.7pt] ($(g0)!\t!(g1)+(0,-2.1mm)$) circle (0.75mm);}
\let\entete\entetecouleur

% --- Encadrés : icône de branche dans l'onglet, un commit du rail sur le bord haut ---
\newlength\railx  \setlength\railx{5mm}   % rail de branche : distance au bord droit
\renewcommand\iconebloc[1]{{\color{cat-#1}\icone{F418}}}
% fer à gauche souple : pas de grands blancs autour du code en ligne
\renewcommand\stylebloc{\rightskip=0pt plus 1fil\relax}
\renewcommand\decorbloc{%
  \filldraw[fill=white,draw=rail,line width=1.3pt] ([xshift=-\railx]frame.north east) circle (1.25mm);}

% --- Colonnes : colonne[branche][couleur] ---
% un rail vertical de la couleur donnée (git-orange par défaut) court le long du bord
% droit (visible entre les blocs), l'étiquette de branche est posée en haut.
\newcommand\rail[1]{\hbox to 0pt{\begin{tikzpicture}[overlay]
    \draw[rail, line width=1.3pt] (\colw-\railx,-0.5\tabh) -- (\colw-\railx,-\colh+2mm);
    \ifx\relax#1\relax\else
      \node[ref=rail, anchor=east, font=\ttfamily\bfseries\fontsize{5.6}{5.6}\selectfont,
            inner ysep=0.8pt] at (\colw,-0.45\tabh) {\icone{F418}\,#1};
    \fi
  \end{tikzpicture}\hss}\nointerlineskip}
\renewcommand\debutcolonne[2]{\ifx\relax#2\relax\else\colorlet{rail}{#2}\fi\rail{#1}}

% --- Mini-graphes de commits ---
% \begin{graphe} ... \end{graphe} : tikzpicture centrée, unités x = 5 mm, y = 3,6 mm
%   \node[c=couleur] (a) at (0,0) {A};      commit (rond, lettre dedans)
%   \node[cm=couleur] ...                    commit de fusion (plein)
%   \node[cg] ...                            commit fantôme (ancien, pointillé)
%   \draw[b=couleur] (a) -- (b);             trait de branche ; (a) |- (b) pour bifurquer
%   \node[ref=couleur, above=.6mm of b] {main};   étiquette de branche · head = HEAD · tag
%   \node[lab] at (..) {texte};              légende en petit · hl = idem en chasse fixe
%   \node[cmd] {git ...};                    commande · fl = flèche · zone = zone Git
\tikzset{
  b/.style={draw=#1, line width=1.2pt, rounded corners=1.4mm},
  c/.style={circle, draw=#1, fill=white, line width=1.2pt, minimum size=3.5mm, inner sep=0pt,
    text=ink, font=\ttfamily\bfseries\fontsize{5.2}{5.2}\selectfont},
  cm/.style={c=#1, fill=#1, text=white},
  cg/.style={c=gutter, dashed, line width=0.8pt, text=gutter},
  ref/.style={fill=#1, text=white, rounded corners=1pt, inner xsep=1.2pt, inner ysep=0.9pt,
    font=\ttfamily\bfseries\fontsize{5}{5}\selectfont},
  head/.style={ref=ink},
  tag/.style={ref=mk-yellow},
  lab/.style={text=ink!75, font=\fontsize{5.6}{6.4}\selectfont, align=center, inner sep=0.5pt},
  hl/.style={lab, font=\ttfamily\fontsize{5.2}{5.2}\selectfont},
  cmd/.style={text=mk-pink, font=\ttfamily\bfseries\fontsize{5.4}{5.4}\selectfont, inner sep=0.8pt},
  fl/.style={-{Stealth[length=1.8mm]}, line width=1pt, draw=ink!70},
  zone/.style={draw=trait, fill=paper, line width=0.8pt, rounded corners=1mm, minimum height=6.5mm,
    text width=13.2mm, align=center, inner sep=0.5pt, font=\bfseries\fontsize{6}{6.8}\selectfont},
}
\newenvironment{graphe}[1][]{\par\vspace{1.2mm}\centering\begin{tikzpicture}[x=5mm,y=3.6mm,#1]}%
  {\end{tikzpicture}\par\vspace{1mm}}

% marques des tableaux comparatifs
\newcommand\oui{\leavevmode{\color{mk-green}\icone{F00C}}}
\newcommand\non{\leavevmode{\color{gutter}--}}
